% Source: https://artofproblemsolving.com/community/c3h1777751_mathematics
% Stable identity: aops-topic-1777751-q01
\textbf{Substitution.} Put $a=\sqrt{\dfrac{yz}{x}}$, $b=\sqrt{\dfrac{zx}{y}}$, $c=\sqrt{\dfrac{xy}{z}}$. Then $bc=x$, $ca=y$, $ab=z$, so $ab+bc+ca=1$, and conversely positive $a,b,c$ with $ab+bc+ca=1$ give positive $x,y,z$ with sum $1$. Dividing numerator and denominator by $x$, $y$, $z$ respectively,
\[
P=\dfrac{1}{1+a^2}+\dfrac{1}{1+b^2}+\dfrac{c}{1+c^2}.
\]
\textbf{Angles.} Write $a=\tan\dfrac {A}{2}$, $b=\tan\dfrac {B}{2}$ with $A,B\in(0,\pi)$. Since $c=\dfrac{1-ab}{a+b}>0$ we have $ab<1$, so $\dfrac{A+B}{2}<\dfrac{\pi}{2}$, and $c=\cot\dfrac{A+B}{2}=\tan\dfrac {C}{2}$ with $C=\pi-A-B>0$. Using $\dfrac{1}{1+\tan^2u}=\cos^2u=\dfrac{1+\cos2u}{2}$ and $\dfrac{\tan u}{1+\tan^2u}=\sin u\cos u=\dfrac{\sin 2u}{2}$,
\[
P=1+\dfrac{\cos A+\cos B}{2}+\dfrac{\sin C}{2}.
\]
\textbf{Bound.} $\cos A+\cos B=2\cos\dfrac{A+B}{2}\cos\dfrac{A-B}{2}\le2\sin\dfrac {C}{2}$, with equality when $A=B$. With $t=\dfrac {C}{2}\in\left(0,\dfrac{\pi}{2}\right)$,
\[
P\le1+\sin t+\sin t\cos t=1+g(t),\qquad g(t)=\sin t\,(1+\cos t).
\]
Now $g'(t)=\cos t+\cos^2t-\sin^2t=(2\cos t-1)(\cos t+1)$, which is positive for $t<\dfrac{\pi}{3}$ and negative after, so $g(t)\le g\left(\dfrac{\pi}{3}\right)=\dfrac{\sqrt3}{2}\cdot\dfrac{3}{2}=\dfrac{3\sqrt3}{4}$.

\textbf{Equality} needs $C=\dfrac{2\pi}{3}$, $A=B=\dfrac{\pi}{6}$: $a=b=\tan15^\circ=2-\sqrt3$, $c=\tan60^\circ=\sqrt3$, i.e.\ $x=y=bc=2\sqrt3-3$ and $z=ab=7-4\sqrt3$ (indeed $x+y+z=1$). Hence the maximum is
\[
\boxed{1+\dfrac{3\sqrt3}{4}}.
\]
